% Answers to the practice exam of Chapter 19, from exam/practice_exam_solutions.md. The evaluation
% questions of lectures/L19/README.md are reflective and have no answers here.

\facitavsnitt{ot}{Övningstentamen}{Hela kursen}
\begin{facit}

\svar{1} $R \approx 0{,}9\,\text{k}\Omega$

\svar{2} $c(c^2 + 4)$

\svar{3} $x = -3$, $y = 3$

\svar{4} \sv{a} $k = -1{,}5$, $m = -1{,}5$ \sv{b} $v = 240^{\circ}$

\svar{5} $x_1 = 3$, $x_2 = -\frac{11}{3}$

\svar{6} \sv{a} Pilar från origo till $(3; 5)$ och till $(2; -4)$.
\sv{b} $\abs{u} = \sqrt{34} \approx 5{,}83$, $\abs{v} = \sqrt{20} \approx 4{,}47$
\sv{c} $v_u \approx 59{,}0^{\circ}$, $v_v \approx -63{,}4^{\circ}$ (eller $296{,}6^{\circ}$)
\sv{d} $w = (-1; 13)$

\svar{7} \sv{a} $f(x) = 7{,}5 \cdot 10^6 \cdot 1{,}025^x$
\sv{b} Cirka $8{,}49 \cdot 10^6\,\text{kr}$
\sv{c} $0 \leq x \leq 30$ år och $7{,}5 \cdot 10^6\,\text{kr} \leq f(x) \leq 15{,}7 \cdot
10^6\,\text{kr}$
\sv{d} Efter cirka $28{,}1$ år.

\svar{8} $f(x) = \frac{3(x - 4)}{x - 3}$, där $x \neq \pm 3$

\svar{9} \sv{a} $u(3) \approx 6{,}43\,\text{V}$
\sv{b} $0 \leq t \leq 10\,\text{s}$ och $2{,}30\,\text{V} \leq u(t) \leq 10\,\text{V}$
\sv{c} Efter cirka $6{,}23\,\text{s}$.

\svar{10} \sv{a} $u(t) = 5\sin\bigl(200\pi t + \frac{\pi}{2}\bigr)\,\text{V}$
\sv{b} Periodtiden är $10\,\text{ms}$. Kurvan börjar i $5\,\text{V}$, är noll vid $2{,}5$ och
$7{,}5\,\text{ms}$ och når $-5\,\text{V}$ vid $5\,\text{ms}$.

\svar{11} $10^{2{,}7} \approx 501$, alltså cirka $500$ gånger.

\svar{12} \sv{a} $f'(x) = 2x - 4$ \sv{b} $x = 2$
\sv{c} Minimum, eftersom $f''(2) = 2 > 0$.
\sv{d} Minsta värdet är $f(2) = -1$.
\sv{e} En uppåtöppen parabel med minimipunkten $(2; -1)$. Den skär $y$-axeln i $(0; 3)$ och
$x$-axeln i $(1; 0)$ och $(3; 0)$, och slutar i $(5; 8)$.

\svar{13} \sv{a} $u'(t) = 5e^{-0{,}5t}\,\text{V/s}$
\sv{b} $u'(2) \approx 1{,}84\,\text{V/s}$
\sv{c} Spänningsökningen minskar: $u''(2) \approx -0{,}92\,\text{V/s}^2 < 0$.

\svar{14} \sv{a} $f'(x) = 2e^{2x}\bigl(2\ln\abs{x} + \frac{1}{x}\bigr)$
\sv{b} $f'(x) = \frac{3x(2\ln x - 1)}{(\ln x)^2}$
\sv{c} $u'(t) = -500\pi\sin\bigl(100\pi t + \frac{\pi}{4}\bigr)\,\text{V/s}$

\svar{15} \sv{a} $I(t) = -5e^{-0{,}1t} + C$
\sv{b} $q(t) = 5 - 5e^{-0{,}1t}$
\sv{c} $q_{\text{tot}}(t) = 6 - 5e^{-0{,}1t}$
\sv{d} Cirka $2{,}65\,\text{C}$

\svar{16} $I = -0{,}25 - j1{,}75\,\text{A} \approx 1{,}8\angle{-1{,}7}\,\text{A}$, med vinkeln i
radianer ($-98{,}1^{\circ}$).

\svar{17} \sv{a} En pil från origo till $10 - j5$, i fjärde kvadranten.
\sv{b} $I \approx 11{,}2e^{-j0{,}46}\,\text{mA}$: $\abs{I} = \sqrt{125} \approx 11{,}2\,\text{mA}$
och $\delta \approx -0{,}46\,\text{rad}$ ($-26{,}6^{\circ}$)
\sv{c} $w = 20\pi \approx 62{,}8\,\text{rad/s}$
\sv{d} $i(t) \approx 11{,}2e^{j(20\pi t - 0{,}46)}\,\text{mA}$

\svar{18} \sv{a} $a = 2 + j$, $b = 3 - j4$, $c = -1 + j3$
\sv{b} Pilar från origo: $a$ i första, $b$ i fjärde och $c$ i andra kvadranten.
\sv{c} $\abs{a} = \sqrt{5} \approx 2{,}24$, $\abs{b} = 5$, $\abs{c} = \sqrt{10} \approx 3{,}16$
\sv{d} $\abs{7 - j7} = 7\sqrt{2} \approx 9{,}9$
\sv{e} $\delta_a \approx 26{,}6^{\circ}$ ($0{,}46\,\text{rad}$), $\delta_b \approx -53{,}1^{\circ}$
($-0{,}93\,\text{rad}$), $\delta_c \approx 108{,}4^{\circ}$ ($1{,}89\,\text{rad}$)
\sv{f} $d \approx -6{,}26 - j3{,}13$

\svar{19} \sv{a} $U_1 \approx 1{,}81 + j0{,}85\,\text{V}$, $U_2 \approx 3{,}54 - j3{,}54\,\text{V}$,
$U_3 \approx 1{,}21 + j0{,}88\,\text{V}$
\sv{b} $U_{\text{tot}} \approx 6{,}56 - j1{,}81\,\text{V}$
\sv{c} $U_1$ och $U_3$ i första kvadranten, $U_2$ och $U_{\text{tot}}$ i fjärde; $U_{\text{tot}}$
ligger strax under den reella axeln.
\sv{d} $u_{\text{tot}}(t) \approx 6{,}81\sin(wt - 15{,}4^{\circ})\,\text{V}$

\end{facit}
